\documentclass[12pt, a4paper]{article}
\usepackage[utf8]{inputenc}
\usepackage[T1]{fontenc}
\usepackage[romanian]{babel}
\usepackage{amsmath}
\usepackage{amssymb}
\usepackage{geometry}
\geometry{a4paper, margin=2cm}
\usepackage{listings}
\usepackage{xcolor}

\lstset{
    language=Python,
    basicstyle=\ttfamily\small,
    keywordstyle=\color{blue},
    stringstyle=\color{red},
    commentstyle=\color{green!50!black},
    showstringspaces=false,
    breaklines=true,
    frame=single,
    numbers=left,
    numberstyle=\tiny\color{gray}
}

\title{Rezolvare Teme - Algoritmi de Primalitate \c{s}i Factorizare}
\author{}
\date{}

\begin{document}
\maketitle
\section*{Tema 5 (Subpunctul 43): Testul Fermat pentru $n = 503$}

\textbf{Cerin\c{t}\u{a}}: Verificarea primalit\u{a}\c{t}ii lui $503$ folosind algoritmul lui Fermat (cel mult 3 martori).

Pentru a fi prim, $n$ trebuie s\u{a} verifice rela\c{t}ia: $b^{n-1} \equiv 1 \pmod n$. \^{I}n cazul nostru, $n-1 = 502$.

\begin{itemize}
    \item \textbf{Martorul $b = 2$}: \\
    Calcul\u{a}m $2^{502} \pmod{503}$. \\
    Folosind ridicarea succesiv\u{a} la putere: $2^9 = 512 \equiv 9 \pmod{503}$. \\
    $2^{502} = (2^9)^{55} \cdot 2^7 \equiv 9^{55} \cdot 128 \pmod{503}$. \\
    Continu\^{a}nd calculele modulare, se ob\c{t}ine $2^{502} \equiv 1 \pmod{503}$.
    
    \item \textbf{Martorul $b = 3$}: \\
    Calcul\u{a}m $3^{502} \pmod{503}$. \\
    De asemenea, se ajunge la rezultatul $3^{502} \equiv 1 \pmod{503}$.
\end{itemize}

\textbf{Concluzie}: Deoarece $b^{n-1} \equiv 1 \pmod n$ este adev\u{a}rat pentru martorii ale\c{s}i, deducem c\u{a} num\u{a}rul $503$ este probabil prim (de fapt, este chiar prim).


\end{document}